
\section{Smoothing Assembly}

We now assemble the stationary martingale bound and the composite remainder via
the Bobkov--Götze smoothing inequality.
% Inputs: the Poisson decomposition, bracket concentration, martingale
% Berry--Esseen bound, and Levin depth-two misadjustment estimate. Together
% they yield a Kolmogorov bound for the stationary \(n_0 = 0\) augmented-chain RR
% statistic.

\textbf{Stationary assembled statistic.} Define
\begin{equation}
\begin{aligned}
S_{n, \text{stat}}^{\text{RR}}(u)
  \coloneqq -\frac{u^\top M_n^{\text{RR}}}{\sqrt{n}} + u^\top \mathcal{R}_{n, \text{stat}}^{\text{RR}},
\end{aligned}
\label{eq:full-decomp}
\end{equation}
\begin{equation}
\begin{aligned}
\mathcal{R}_{n, \text{stat}}^{\text{RR}}
  \coloneqq D_{2, n}^{\text{RR}} + R_n^{\text{mis, RR}},
\end{aligned}
\end{equation}
where both terms are read under the stationary augmented-chain convention.
Dividing by \(\sigma_n^{\text{RR}}(u)\),
% The deterministic-start transient and random initial-product discrepancy are
% not part of this stationary result; they are handled in the burn-in transfer
% theorem.
\begin{equation}
\begin{aligned}
\frac{S_{n, \text{stat}}^{\text{RR}}(u)}{\sigma_n^{\text{RR}}(u)}
  = X_n + Y_n,
\end{aligned}
\label{eq:XY-split}
\end{equation}
\begin{equation}
\begin{aligned}
X_n \coloneqq -\frac{u^\top M_n^{\text{RR}}}{\sqrt{n} \, \sigma_n^{\text{RR}}(u)},
\quad
Y_n \coloneqq \frac{u^\top \mathcal{R}_{n, \text{stat}}^{\text{RR}}}{\sigma_n^{\text{RR}}(u)}.
\end{aligned}
\end{equation}
The martingale bound applies to \(X_n\) with direction \(-u\).
% The bounded-increment constant and predictable variance are unchanged:
% \(\varkappa(-u) = \varkappa(u)\) and
% \(\sigma_n^{2, \text{RR}}(-u) = \sigma_n^{2, \text{RR}}(u)\).

\textbf{Smoothing inequality.} For any \(t > 0\),
\begin{equation}
\begin{aligned}
d_K (X + Y, \mathcal{N}(0, 1))
  \le d_K (X, \mathcal{N}(0, 1))
   + \mathbb{P}(|Y| > t)
   + \frac{t}{\sqrt{2 \pi}},
\end{aligned}
\label{eq:smoothing}
\end{equation}
Markov's inequality with \(t = e \, \lVert Y \rVert_{L_p}\) gives
\begin{equation}
\begin{aligned}
d_K (X + Y, \mathcal{N}(0, 1))
  \le d_K (X, \mathcal{N}(0, 1))
   + \frac{e \, \lVert  Y  \rVert_{L_p}}{\sqrt{2 \pi}}
   + e^{-p},
\end{aligned}
\label{eq:smoothing-Lp}
\end{equation}
% The third term uses the Lipschitz constant of the standard normal cdf. The
% trailing \(e^{-p}\) is \(O(1 / n)\) when \(p \ge \log n\).

\textbf{\(L_p\)-bound on the composite remainder.} The two stationary remainder pieces
contribute additively.
% The finite-start deterministic transient is deliberately excluded: setting
% \(\theta_0 = \theta^*\) removes only that deterministic term, not the startup
% discrepancy between zero-start perturbation variables and the stationary
% augmented chain.

\begin{lemma}\label{lem:R-bound}
  \textbf{(Stationary composite remainder bound.)}
  Assume the hypotheses and step-size restrictions of the stationary
  PR-averaged RR misadjustment bound above. Then for every \(u \in \mathbb{R}^d\),
  every \(p \ge 2\), and every \(q \ge 2\) satisfying \(p \le q / 2\) and
  \(2 \alpha \le \alpha_{\text{stat}}(q)\),
\begin{equation}
\begin{aligned}
  \lVert  u^\top \mathcal{R}_{n, \text{stat}}^{\text{RR}}  \rVert_{L_p}
    &\le \frac{C_{\text{D2}} \, \lVert  u  \rVert}{\sqrt{n}}
     + \lVert  u  \rVert \, C_{\text{mis}} \, \Biggl(
        \sqrt{n} \, \alpha^2
        + (1 + d^{1 / q}) \, p^{7 / 2} \, t_{\text{mix}}^{5 / 2}
            \, \sqrt{n} \, \alpha^{3 / 2} \, \log^{3 / 2}(1 / (\alpha a))
      \\
        & \quad \quad \quad \quad
        + p^{3 / 2} \, \sqrt{\alpha}
        + p^3 \, (\alpha n)^{-1 / 2} \, \log^{1 / 2}(1 / (\alpha a))
        + \Phi_+(p, \alpha) \, n^{-1 / 2}
      \Biggr),
\end{aligned}
\end{equation}
  with constants
\begin{equation}
\begin{aligned}
  C_{\text{D2}} \coloneqq 3 \, t_{\text{mix}} \, \lVert  \varepsilon  \rVert_\infty
              \, (C_{\mathcal{Q}} + C_2 / a^2),
\end{aligned}
\end{equation}
  and \(C_{\text{mis}}\) the constant in the misadjustment theorem.
\end{lemma}

\textit{Proof.} Combine the sup-norm bound for \(D_{2,n}^{\text{RR}}\) with
Theorem~\ref{thm:misadjustment}. \(\square\)

\begin{theorem}\label{thm:RR-BE}
  \textbf{(Stationary augmented-chain Berry--Esseen assembly for the RR comparison statistic.)}
  Assume \textbf{UGE 1}, \(\pi(\widetilde{A}) = 0\), \(\pi(\varepsilon) = 0\),
  \(\lVert  \varepsilon  \rVert_\infty < \infty\), \(\sigma^2(u) > 0\),
  \(0 < \alpha\), and the stationary augmented-chain
  convention above. Use the external inputs summarized in
  Section~\ref{sec:external-inputs}. Set
  \(p = \max(2, \left\lceil \log n \right\rceil)\) and \(q = \max(2 p, \left\lceil \log(e \, d) \right\rceil, 2)\).
  Then \(p \ge 2\), \(q \ge 2\), and \(p \le q / 2\). If \(n \ge 3\) satisfies the
  small-step and variance lower-bound conditions
\begin{equation}
\begin{aligned}
  2 \alpha \le \alpha_{\text{stat}}(q),
  \quad
  n \, \alpha \, a \ge \frac{2 \, C_3 \, \lVert  u  \rVert^2}{\sigma^2(u)},
\end{aligned}
\label{eq:RR-BE-master-conditions}
\end{equation}
  where the second inequality is \eqref{eq:variance-lb-condition}, then
\begin{equation}
\begin{aligned}
  d_K \left(
    \frac{S_{n, \text{stat}}^{\text{RR}}(u)}{\sigma_n^{\text{RR}}(u)},
    \mathcal{N}(0, 1)
  \right)
    &\le \frac{C_{K, 1}(u) \, \log^{3 / 4} n}{n^{1 / 4}}
     + \frac{C_{K, 2}(u) \, \log n}{\sqrt{n}} \\
    &\quad + \frac{e \, \lVert  u^\top \mathcal{R}_{n, \text{stat}}^{\text{RR}}  \rVert_{L_p}}{\sqrt{2 \pi} \, \sigma_n^{\text{RR}}(u)}
     + \frac{e}{n},
\end{aligned}
\label{eq:RR-BE-master}
\end{equation}
  with \(\lVert  u^\top \mathcal{R}_{n, \text{stat}}^{\text{RR}}  \rVert_{L_p}\) bounded by the preceding
  lemma, and \(C_{K, 1}(u), C_{K, 2}(u)\) the constants of the martingale
  Berry--Esseen theorem.
\end{theorem}

\textit{Proof.} Apply \eqref{eq:smoothing-Lp} with \(X = X_n\) and \(Y = Y_n\). The martingale
terms come from Theorem~\ref{thm:M-RR-BE}, and the remainder term is
\[
\frac{e\lVert Y_n\rVert_{L_p}}{\sqrt{2\pi}}
  =
\frac{e\lVert u^\top\mathcal{R}_{n,\mathrm{stat}}^{\mathrm{RR}}\rVert_{L_p}}
     {\sqrt{2\pi}\,\sigma_n^{\mathrm{RR}}(u)}.
\]
Since \(p\ge\log n\),
\(e^{-p} \le n^{-1}\). \(\square\)

% \begin{corollary}\label{cor:RR-BE-admissible-alpha}
% \textbf{(Stationary triangular-array admissible-step CLT.)}
% Assume the stationary augmented-chain hypotheses of Theorem~\ref{thm:RR-BE}, including
% \(\sigma^2(u) > 0\). Let \(\alpha = \alpha_n\) be a step-size sequence and set
% \(p_n \coloneqq \max(2, \left\lceil \log n \right\rceil)\), \(q_n \coloneqq \max(2 p_n, \left\lceil \log(e \, d) \right\rceil, 2)\),
% and \(\Lambda_n \coloneqq \log(1 / (\alpha_n a))\). Then \(p_n \le q_n / 2\).
% Assume that, for all sufficiently large \(n\),
% \begin{equation}
% \begin{aligned}
%   2 \alpha_n \le \alpha_{\text{stat}}(q_n),
%   \quad
%   n \, \alpha_n \, a \ge \frac{2 \, C_3 \, \lVert  u  \rVert^2}{\sigma^2(u)}.
% \end{aligned}
% \end{equation}
% If
% \begin{equation}
% \begin{aligned}
%   p_n^3 (n \alpha_n)^{-1 / 2} \Lambda_n^{1 / p_n} \to 0,
%   \quad
%   p_n^{7 / 2} \sqrt{n} \, \alpha_n^{3 / 2} \Lambda_n^{3 / 2} \to 0,
% \end{aligned}
% \label{eq:alpha-window}
% \end{equation}
% then
% \begin{equation}
% \begin{aligned}
%   d_K \left(
%     \frac{S_{n, \text{stat}}^{\text{RR}}(u)}{\sigma_n^{\text{RR}}(u)},
%     \mathcal{N}(0, 1)
%   \right) \to 0.
% \end{aligned}
% \end{equation}
% In particular, for a power scale \(\alpha_n = c \, n^{-\gamma}\), up to
% logarithmic factors, the admissible window is
% \begin{equation}
% \begin{aligned}
% \frac{1}{3} < \gamma < 1.
% \end{aligned}
% \end{equation}
% \end{corollary}

% \textit{Proof.} Apply Theorem~\ref{thm:RR-BE} and Lemma~\ref{lem:R-bound} with
% \(p = p_n\), \(q = q_n\), and \(\alpha = \alpha_n\). The displayed conditions are the
% remaining nontrivial remainder constraints. For \(\alpha_n = c n^{-\gamma}\) they
% reduce, modulo logarithms, to
% \(n^{-(1 - \gamma) / 2} \to 0\) and
% \(n^{1 / 2 - 3 \gamma / 2} \to 0\), i.e. \(\gamma < 1\) and
% \(\gamma > 1 / 3\). \(\square\)

\begin{corollary}\label{cor:RR-BE-working}
  \textbf{(Stationary balanced-scale augmented-chain Berry--Esseen bound.)}
  At the balanced scale \(\alpha = c \, n^{-1 / 2}\) with \(c > 0\), put
  \(p = \max(2, \left\lceil \log n \right\rceil)\) and \(q = \max(2 p, \left\lceil \log(e \, d) \right\rceil, 2)\).
  Then \(p \le q / 2\). Assume the stationary augmented-chain hypotheses of
  Theorem~\ref{thm:RR-BE}, including \(\sigma^2(u) > 0\). If \(n \ge 3\) satisfies
\begin{equation}
\begin{aligned}
  2 \alpha \le \alpha_{\text{stat}}(q),
  \quad
  n \, \alpha \, a \ge \frac{2 \, C_3 \, \lVert  u  \rVert^2}{\sigma^2(u)},
\end{aligned}
\end{equation}
  the bound
  \eqref{eq:RR-BE-master} reduces to
\begin{equation}
\begin{aligned}
  d_K \left(
    \frac{S_{n, \text{stat}}^{\text{RR}}(u)}{\sigma_n^{\text{RR}}(u)},
    \mathcal{N}(0, 1)
  \right)
    \le \frac{C(u) \, \text{polylog}(n)}{n^{1 / 4}},
\end{aligned}
\end{equation}
  where \(C(u)\) depends on \(\lVert  u  \rVert\), \(\sigma(u)\), \(C_{\mathcal{Q}}\),
  \(\lVert  \overline{A}  \rVert\), \(\lVert  \overline{A}^{-1}  \rVert\), \(\kappa_Q\), \(C_A\),
  \(\lVert  \varepsilon  \rVert_\infty\), \(\lVert  \Sigma_{\varepsilon}^{(\text{M})}  \rVert\),
  \(t_{\text{mix}}\), \(a\), \(\alpha_\infty\), \(c\), and the universal constants of
  the smoothing inequality, Lemma~\ref{lem:external-martingale-be}, and
  \eqref{eq:external-markov-conc}.
\end{corollary}

% \textit{Proof.} At \(\alpha = c \, n^{-1 / 2}\),
% \(C_{\text{D2}} \lVert  u  \rVert / \sqrt{n} = O(n^{-1 / 2})\);
% \(\sqrt{n} \alpha^2 = c^2 \, n^{-1 / 2}\);
% \(p^{7/2}\sqrt{n}\,\alpha^{3/2}\log^{3/2}(1/(\alpha a))
% =O(\operatorname{polylog}(n)n^{-1/4})\);
% \(p^{3 / 2} \sqrt{\alpha} = O(\text{polylog}(n) n^{-1 / 4})\);
% \(p^3(\alpha n)^{-1/2}\log^{1/2}(1/(\alpha a))
% =O(\operatorname{polylog}(n)n^{-1/4})\); and
% \(\Phi_+(p, \alpha) n^{-1 / 2} = O(\text{polylog}(n) n^{-1 / 2})\).
% After division by \(\sigma_n^{\text{RR}}(u) \ge \sigma(u) / \sqrt{2}\), all remainder
% terms are dominated by \(\text{polylog}(n) n^{-1 / 4}\). Combine with the
% martingale Berry--Esseen term.
% \(\square\)

% \begin{corollary}\label{cor:RR-BE-sigma}
% \textbf{(Stationary asymptotic-normalization version.)}
% Assume the stationary augmented-chain hypotheses of Theorem~\ref{thm:RR-BE} for the same
% \(n,p,q,\alpha,u\). In particular, assume \(p \le q / 2\), \(\sigma^2(u) > 0\),
% and
% \begin{equation}
% \begin{aligned}
%   2 \alpha \le \alpha_{\text{stat}}(q),
%   \quad
%   n \, \alpha \, a \ge \frac{2 \, C_3 \, \lVert  u  \rVert^2}{\sigma^2(u)}.
% \end{aligned}
% \end{equation}
% Then the same finite-\(n\) bound with asymptotic normalization contains one
% additional variance-comparison term:
% \begin{equation}
% \begin{aligned}
%   d_K \left(
%     \frac{S_{n, \text{stat}}^{\text{RR}}(u)}{\sigma(u)},
%     \mathcal{N}(0, 1)
%   \right)
%     &\le \frac{C_{K, 1}(u) \, \log^{3 / 4} n}{n^{1 / 4}}
%      + \frac{C_{K, 2}(u) \, \log n}{\sqrt{n}} \\
%     &\quad + \frac{e \, \lVert  u^\top \mathcal{R}_{n, \text{stat}}^{\text{RR}}  \rVert_{L_p}}{\sqrt{2 \pi} \, \sigma_n^{\text{RR}}(u)}
%      + \frac{e}{n}
%      + \frac{C \, \lVert  u  \rVert^2}{n \, \alpha \, a \, \sigma^2(u)}.
% \end{aligned}
% \label{eq:RR-BE-sigma-master}
% \end{equation}
% Consequently, under the balanced-scale hypotheses above,
% \begin{equation}
% \begin{aligned}
%   d_K \left(
%     \frac{S_{n, \text{stat}}^{\text{RR}}(u)}{\sigma(u)},
%     \mathcal{N}(0, 1)
%   \right)
%     \le \frac{C'(u) \, \text{polylog}(n)}{n^{1 / 4}}.
% \end{aligned}
% \end{equation}
% \end{corollary}

% \textit{Proof.} Set \(r \coloneqq \sigma_n^{\text{RR}}(u) / \sigma(u)\) and write
% \(W \coloneqq S_{n, \text{stat}}^{\text{RR}}(u) / \sigma_n^{\text{RR}}(u)\), so
% \(W r = S_{n, \text{stat}}^{\text{RR}}(u) / \sigma(u)\). Under
% the variance lower-bound condition \eqref{eq:variance-lb-condition}, \(r \in [1 / \sqrt{2}, \, r_{\max}]\) for some
% finite \(r_{\max}\) (from the trivial upper bound
% \(\sigma_n^{2, \text{RR}}(u) \le C_{\mathcal{Q}}^2 \, \lVert  \Sigma  \rVert \, \lVert  u  \rVert^2\)).
% For \(r\) in this compact interval,
% \begin{equation}
% \begin{aligned}
% \sup_x | \Phi(x / r) - \Phi(x) | \le C_\Phi \, |r - 1|,
% \quad
% C_\Phi \coloneqq \sqrt{2} / \sqrt{\pi e},
% \end{aligned}
% \end{equation}
% Hence
% The displayed Lipschitz constant follows from
% \(\sup_x |x \, \phi(x)| = 1 / \sqrt{2 \pi e}\).
% \begin{equation}
% \begin{aligned}
% d_K (W r, \mathcal{N}(0, 1))
%   \le d_K (W, \mathcal{N}(0, 1)) + C_\Phi \, |r - 1|.
% \end{aligned}
% \end{equation}
% The variance comparison of Section 4.5 gives
% \begin{equation}
% \begin{aligned}
% |r - 1|
%   = \frac{|\sigma_n^{2, \text{RR}}(u) - \sigma^2(u)|}{\sigma(u) \, (\sigma_n^{\text{RR}}(u) + \sigma(u))}
%   \le \frac{C_3 \, \lVert  u  \rVert^2}{n \, \alpha \, a \, \sigma^2(u)},
% \end{aligned}
% \end{equation}
% using \(\sigma_n^{\text{RR}}(u) + \sigma(u) \ge \sigma(u)\). At
% \(\alpha = c n^{-1 / 2}\) the final term is \(O(n^{-1 / 2})\). \(\square\)

% This is an \(n_0 = 0\) stationary augmented-chain statement; deterministic starts
% are transferred in the burn-in chapter.
